\section{Introduction}\label{sec:introduction}

Let $p$ be a prime, let $k/\mathbb Q_p$ be a finite extension, and let
$X/k$ be a smooth proper geometrically connected curve of genus at least
two. Choose an algebraic closure $\bar k$ and a geometric base point
of $X$. The full arithmetic \'etale fundamental group fits into
\begin{equation}\label{eq:fundamental-sequence}
 1\longrightarrow\Delta_X\longrightarrow\Pi_X
   \longrightarrow G_k\longrightarrow1,
 \qquad
 \Delta_X=\pi_1(X_{\bar k}),\quad \Pi_X=\pi_1(X),\quad
 G_k=\Gal(\bar k/k).
\end{equation}
A rational point $a\in X(k)$ induces a continuous group-homomorphic
section $s_a:G_k\to\Pi_X$. Its class under conjugation by $\Delta_X$
is independent of the path used to identify base points. The local
$p$-adic section conjecture asserts that this construction accounts for
all sections, and distinguishes all rational points.

A section whose class comes from a rational point is called
\emph{geometric}. The conjecture thus asks whether rational points can
be recovered from the splittings of the arithmetic fundamental group.
It is a local analogue of the section conjecture in Grothendieck's
letter to Faltings of 27 June 1983, formulated over fields finitely
generated over the prime field
\cite[equation~(7) and the subsequent properness correction]{GrothendieckLetter}.
The local question retains the full arithmetic and geometric groups
in~\eqref{eq:fundamental-sequence}, while bringing the geometry of
$p$-adic valuations and reduction into the problem.

\begin{theorem}\label{thm:section}
For every prime $p$, every finite extension $k/\mathbb Q_p$, and every
smooth proper geometrically connected curve $X/k$ of genus at least two,
the map
\[
 X(k)\longrightarrow
 \left\{\begin{gathered}
 s:G_k\to\pi_1(X)\text{ a continuous group homomorphism}\\
 \text{such that }(\pi_1(X)\to G_k)\circ s=\mathrm{id}_{G_k}
 \end{gathered}\right\}\big/\pi_1(X_{\bar k})
\]
is bijective. The quotient is by conjugation by the full geometric
\'etale fundamental group.
\end{theorem}

\subsection{Earlier results and the lifting problem}

The birational analogue replaces the arithmetic \'etale fundamental
group by the absolute Galois group of the curve's function field.
For a Galois extension with a chosen valuation, the
\emph{decomposition group} consists of the automorphisms preserving
its valuation ring.
Koenigsmann proves that every birational section over a finite extension
of $\mathbb Q_p$ is contained in the decomposition group of a rational
point \cite[preprint version, Proposition~2.4(2)]{Koenigsmann}.
Pop subsequently obtains a birational criterion from sections of an
elementary abelian quotient that lift to the corresponding elementary
metabelian quotient
\cite[Theorems~A and~A$'$]{PopBirational}.
These results make the passage from an \'etale section to a birational
section a natural objective. Such a lift is additional information:
a section of the quotient group in~\eqref{eq:fundamental-sequence} does
not by itself supply a splitting of the absolute Galois sequence.

Pop--Stix localize every \'etale section at a valuation extending the
$p$-adic valuation \cite[preprint version, Theorem~26]{PopStix}.
This reduces the problem to understanding which valuations can support
a section. Pajwani develops this valuative approach through \'etale
homotopy and fixed points in Berkovich spaces
\cite[preprint version, Theorem~1.4 and Corollary~1.5]{Pajwani}.
The passage from a fixed valuation to a rational point remains the
geometric issue addressed here.

Recent results also constrain the possible nongeometric sections.
Bresciani introduces toric sections and proves the section conjecture
for them; they map
injectively to the set of \'etale sections
\cite[Definition~3.7 and Theorems~1.4 and~5.10]{BrescianiToric}.
For an arbitrary nongeometric \'etale section, his Theorem~1.5 produces
a finite \'etale neighbourhood whose index is divisible by $p$.
Here a neighbourhood is a connected finite \'etale cover to which the
section lifts, and the index of a curve is the greatest common divisor
of the degrees of its closed points.
In a different direction, Betts--Stix prove finiteness of the local
image of global Selmer sections for smooth proper geometrically
connected curves of genus at least two over number fields containing no CM subfield
\cite[preprint version, Theorem~A]{BettsStix}.
Their Selmer sections already arise from points at every completion;
the finiteness assertion concerns the local images of this globally
constrained set of sections.

Our lifting step is a statement about completions of the universal
\'etale function field. It compares the absolute and \'etale
decomposition groups at every rank-one valuation extending the
$p$-adic norm. This supplies the birational lift at precisely the
valuations that must be excluded from the Pop--Stix localization.

\subsection{The completion statement}

Put $K=k(X)$ and choose a separable closure $K^{\mathrm{sep}}$ containing
$\bar k$. Let $\widetilde K\subset K^{\mathrm{sep}}$ be the compositum
of the function fields of all connected finite \'etale covers of $X$,
with all their $K$-embeddings in $K^{\mathrm{sep}}$. Then
$\Gal(\widetilde K/K)=\Pi_X$ after the chosen base point identification.
A real-valued rank-one norm on $K$ means a nonarchimedean field absolute
value with values in $\mathbb R_{\ge0}$; in particular its support is
zero. We normalize an extension of the $p$-adic norm by $|p|=p^{-1}$.

\begin{theorem}\label{thm:completion}
Let $p$ be a prime, let $k/\mathbb Q_p$ be finite, and let $X/k$ be a
smooth proper geometrically connected curve of genus at least two.
Put $K=k(X)$, choose a separable closure $K^{\mathrm{sep}}$, and let
$\widetilde K\subset K^{\mathrm{sep}}$ be the compositum of the function
fields of all connected finite \'etale covers of $X$, with all their
$K$-embeddings.
For every real-valued rank-one norm on $K$ extending the normalized
$p$-adic norm on $k$, and every prolongation to $K^{\mathrm{sep}}$, set
\[
 E=\widehat{K^{\mathrm{sep}}},\qquad
 H=\widehat{\widetilde K}\subset E,
\]
where the completions use that fixed prolongation.
Then the specified subfield $K^{\mathrm{sep}}$ of $E$ is contained in $H$.
\end{theorem}

The statement concerns the chosen copies of these fields inside one
completion. Write $D_\nu(L/K)$ for the decomposition group of a
valuation $\nu$ on a Galois extension $L/K$.
Corollary~\ref{cor:decomposition} uses Theorem~\ref{thm:completion} to
identify the decomposition group in $K^{\mathrm{sep}}/K$ with the one
in $\widetilde K/K$. Indeed, an automorphism in the kernel fixes a
dense subfield of $H$ and hence fixes $K^{\mathrm{sep}}$.

A section localized at a rank-one valuation therefore lifts to a
birational section. Koenigsmann's theorem places this lifted section
also at a rational point. Its fixed field would then carry two
inequivalent henselian rank-one valuations, one extending the $p$-adic
valuation and the other trivial on constants. The contradiction is an
explicit instance of the classical two-valuation principle associated
with F.~K.~Schmidt \cite{SchmidtValuations}; Section~\ref{sec:sections}
gives the needed approximation and square-root argument, including
$p=2$.

Rank-two valuations are handled by their rank-one coarsenings. The
remaining case is a closed-point valuation refined by the $p$-adic
valuation on its residue field, called type~$2h$ by Pop--Stix.
The surjection of the section onto $G_k$ forces that closed point to
be rational, and their type-$2h$ analysis then identifies the section
\cite[preprint version, Section~6.1]{PopStix}.
Section~\ref{sec:sections} completes this argument and proves
injectivity of the point-to-section map by Kummer theory on the
Jacobian.

\subsection{Two stages of collar roots}

The geometric content lies in Theorem~\ref{thm:completion}.
A finite extension of $\bar k(X)$ is the function field of a finite
curve cover, which may have branch points. Near such a point, with
local coordinate $t$, a branch has the form $t=z^e$.
If another cover supplies a root of $t$ of order divisible by all
these branch indices, the pulled-back cover splits over an annulus
around the branch point. One can then replace the part over the inner
disk by a trivial cover and glue to obtain an \'etale cover.
The roots must exist on the \emph{entire} inverse image of the annulus:
a root on only one sheet does not trivialize the full cover needed
for gluing. We use this operation twice, with different sources of
collar roots.

This gluing and algebraization method belongs to the formal and rigid
analytic patching tradition of Harbater and Liu
\cite{HarbaterPatching,LiuPatching}. Poineau treats the corresponding
Berkovich framework \cite[Section~1 and Appendix~A]{GAGA}.
Section~\ref{sec:patching} proves the precise replacement and
field-embedding statements used here. In particular, evaluating the
resulting exterior analytic map retains the specified valued field
embedding; the map itself need not descend to $\bar k$.

Fix an odd prime $\ell\ne p$. In the first stage, coordinates on the
node annuli of a stable model acquire $\ell$-th roots on suitable
finite \'etale covers. The construction uses the lifting theorem of
Amini--Baker--Brugall\'e--Rabinoff, which realizes tame coverings of
metrized complexes as curve covers \cite[Theorem~7.7]{ABBR}.
A metrized complex records the skeleton together with the residue
curve at each vertex and the marked points corresponding to edge ends.
Small collars around arbitrary branch points need not themselves be
stable node annuli. The resolution of nonsingularities theorem of
Mochizuki--Tsujimura supplies a finite \'etale cover whose stable model
maps to a prescribed semistable model
\cite[Theorem~A and Proposition~2.4(iv)]{MochizukiTsujimura}.
This extends Tamagawa's result for closed points on a stable model
\cite[Theorem~0.2(v)]{TamagawaRNS} and Lepage's corresponding resolution
theorem for Mumford curves \cite[Theorem~2.7]{LepageRNS}.
The general theorem of Mochizuki--Tsujimura allows us to move these
collars into the stable-node construction.
Patching then places in $H$ every finite Galois extension whose
horizontal ramification indices divide $\ell$. Here horizontal
indices refer to ramification at closed points over $\bar k$.
Section~\ref{sec:tame-roots} proves this first stage.

The second stage produces roots of arbitrary positive order $m$ while
keeping the auxiliary cover's horizontal indices dividing
$\ell$. We construct a cyclic cover of the projective line whose
skeleton has two vertices and $\ell$ parallel edges. A cycle using
two distinct edges gives a character of the torus uniformizing its
Jacobian. Along a chosen annulus that character has slope one, so it
is a constant times the annular coordinate times an analytic unit
close to one. Figure~\ref{fig:parallel-edges} shows the geometry.
This uses the nonarchimedean uniformization theory of Raynaud and
Bosch--L\"utkebohmert \cite{RaynaudUniformization,BoschLutkebohmert}
and the compatibility of classical and tropical Abel--Jacobi maps
proved by Baker--Rabinoff
\cite[Proposition~6.1 and Corollary~6.6]{BakerRabinoff}.

Pullback by multiplication by $m$ on the Jacobian supplies a root of
the character on the whole inverse annulus. On a sufficiently long
middle collar, the remaining unit has an $m$-th root by the $p$-adic
logarithm and exponential. This yields a root of the coordinate itself.
The auxiliary covers still have horizontal indices dividing $\ell$,
so the first stage places their function fields in $H$.
Section~\ref{sec:arbitrary-roots} constructs these roots, and a second
application of patching in Section~\ref{sec:realization} realizes
arbitrary finite Galois extensions in the same completion.
The common requirement of roots on every inverse-image sheet is what
allows both stages to use the same patching proposition.

\subsection{Global consequences}

The local theorem also identifies the possible localizations of global
sections. Let $C/F$ be a smooth proper geometrically connected curve of
genus at least two over a number field. A global section determines a
point of $C(F_v)$ at every finite place by Theorem~\ref{thm:section},
and a connected component of $C(F_v)$ at every real place by the real
section theorem \cite[Introduction, Theorem]{BrescianiVistoliReal}.
Together these form a \emph{modified adelic point}; complex places
contribute no factor. The finite-cover descent locus consists of the
tuples that lift to some global twist of each torsor under a finite
$F$-group scheme. Proposition~\ref{prop:localization-image} shows that
this locus is exactly the image of global sections under localization.
Both inclusions use the descent bridge of Harari--Stix
\cite[preprint version, Theorem~2.1 and Remark~2.2(1)]{HarariStixFiniteDescent}.

When this finite-cover locus consists exactly of the diagonal rational
points, Stix's finite-support theorem identifies every global section
with a point section
\cite[arXiv version, Corollary~5]{StixFiniteSupport}.
The resulting bijection is Theorem~\ref{thm:global-descent}.
Stoll's finite-descent criteria supply the required equality when the
curve has a nonconstant morphism over its number field to an abelian
variety with finitely many rational points and trivial divisible
Tate--Shafarevich subgroup, or when its Jacobian has these properties.
They also cover the smooth projective modular curves $X_0(N)$ and
$X_1(N)$ of genus at least two over $\mathbb Q$
\cite[Corollary~8.1, Theorem~8.6, and Corollary~8.8]{StollFiniteDescent},
with the corrections in \cite{StollFiniteDescentErrata}.
The Selmer converse in the companion article
\cite[Theorem~1.1]{OpenAISelmerConverse2026} gives a further arithmetic
test for an elliptic target over $\mathbb Q$: vanishing of its full
$p$-power Selmer corank at one prime.
Section~\ref{sec:global-sections} states the precise hypotheses and
deductions for these cases.

The localization description also yields constraints without assuming
that the finite-cover locus consists of rational points. Every ordinary
adelic representative of a localization tuple survives linear descent
and lies in the Brauer--Manin set, by further results of Harari--Stix
\cite[preprint version, Corollary~3.1 and Theorem~4.1]{HarariStixFiniteDescent}.
If $F$ contains no CM subfield, Betts--Stix's theorem makes the projection
of the localization image to $C(F_v)$ finite at every finite place $v$
\cite[preprint version, Theorem~A]{BettsStix}.
These statements constrain the localization image; the finite-cover
equality remains a separate arithmetic hypothesis in the global
section-conjecture criterion.
